% Generated draft: requires mathematical and source review.
\documentclass[11pt]{article}
\usepackage[margin=25mm]{geometry}
\usepackage{fontspec}
\setmainfont{Latin Modern Roman}
\usepackage{amsmath,amssymb,mathtools,unicode-math}
\setmathfont{Latin Modern Math}
\usepackage{enumitem,needspace,xurl,hyperref}
\setlength{\parindent}{0pt}
\setlength{\parskip}{5pt}
\setlength{\emergencystretch}{3em}
\newcommand{\N}{\mathbb{N}}
\newcommand{\Z}{\mathbb{Z}}
\newcommand{\Q}{\mathbb{Q}}
\newcommand{\R}{\mathbb{R}}
\newcommand{\C}{\mathbb{C}}
\newcommand{\F}{\mathbb{F}}
\newcommand{\A}{\mathbb{A}}
\newcommand{\PP}{\mathbb{P}}
\newcommand{\E}{\mathbb{E}}
\newcommand{\eps}{\varepsilon}
\newcommand{\dd}{\,\mathrm{d}}
\begin{document}
\section*{UnsolvedMath ID: 1200036}
\textbf{Source problem code:} catalogue entry
\subsection{Definitions and mathematical statement}
In the local-weak interpretation, let $X_n$ be finite nonempty simple graphs and choose a vertex $v_n$ uniformly from $V(X_n)$. Write $X_n\to\mu$ when, for every radius $r\geq0$ and finite rooted graph $(B,o)$, the probability that the induced rooted radius-$r$ ball about $v_n$ is isomorphic to $(B,o)$ converges to the corresponding probability under a law $\mu$ on rooted locally finite graphs. A spanning tree of $X_n$ is a connected acyclic subgraph with vertex set $V(X_n)$.

\textbf{Question.} Suppose that
\[
G_n\to\mu,\qquad H_n\to\mu,\qquad T_n\to\nu,
\]
where each $T_n$ is a spanning tree of $G_n$. Must there exist a spanning tree $P_n$ of $H_n$ for every $n$ such that $P_n\to\nu$? Explicitly, must one be able to choose all $P_n$ so that, for every rooted radius-$r$ pattern $(B,o)$,
\[
\mathbb P\big(B_r(P_n,v_n)\cong(B,o)\big)
-\mathbb P\big(B_r(T_n,w_n)\cong(B,o)\big)\longrightarrow0,
\]
where $v_n$ and $w_n$ are uniform vertices of $H_n$ and $G_n$, respectively?

The source does not explicitly require $H_n$ to be connected, although that is necessary for a spanning tree. The substantive connected-graph variant adds that assumption; it is not implied merely by having the same local limit as $G_n$. The source also leaves the graph-limit topology and any uniform degree bound unspecified. The local interpretation above follows the research summary, but does not silently replace other possible notions of graph convergence. Only the limits of the tree graphs themselves are requested, not convergence of jointly marked pairs $(G_n,T_n)$ and $(H_n,P_n)$.
\subsection{Short English statement}
Can a limiting distribution of spanning-tree neighborhoods obtained from one graph sequence also be realized by spanning trees in every other sequence with the same ambient local limit? Connectedness and the graph-convergence convention must be fixed.
\subsection{Sources}
\begin{itemize}
\item M. Ab\'ert, \emph{Some questions} (2010), Question 36, \url{https://www.renyi.hu/~abert/questions.pdf}; manual transcription, accessed 2026-07-29. The background says current status was not independently verified.
\item The supplied research summary mentions a subgraph-lifting-of-limits program but gives no publication, theorem, or author attribution for it or a resolution. Its open assessment is expressly provisional and search-cap limited.
\item Local registry \texttt{data/unsolvedmath/problems.json}, permanent ID $1200036$, source code catalogue entry, \emph{Some Questions --- Question 36}. The record labels this open and OPEN-TRIAGE, with NEEDS\_REVIEW source-status and rights qualifications.
\end{itemize}
\end{document}
